\speciestitle{Sulfur Dioxide}{SO\cs2}{SO2}
\begin{description}
\item[Useful range:] 215\,--\,10.0\,hPa
\item[Contact:] William Read, \email{William.G.Read@jpl.nasa.gov}
\end{description}

\subsection*{Introduction}

The standard SO\cs2 product is taken from the 240-GHz retrieval. MLS can
only measure significantly enhanced concentrations above nominal
background such as that from volcanic injections.  SO\cs2 has not yet
been validated.

\subsection*{Changes from v2.2}

The \thisv SO\cs2 retrieval will be impacted by the addition of interline
interference terms in the O\cs3 line shape model, an updated CO line
width parameter, and using a different set of 240\,GHz channels.
Another \thisv change is the spectral baseline treatment that now uses a
frequency-squared extinction term, configured as a relative
humidity-like (RH) species. The RH treatment responds better to high
extinction conditions when the middle troposphere is very humid. This
enhancement eliminated a high bias in 316\,hPa O\cs3 in the tropics that
is present in v2. An unfortunate side effect of the RH baseline is that
it is more adversely affected by clouds, causing spikes in the retrieval
of R3 products including SO\cs2.

\subsection*{Resolution}
\onestandardavk{SO2}{SO\cs2}%

Based on Figure~\ref{fig:AvkSO2}, the vertical resolution for SO\cs2 is
\texttildelow3\,km and the horizontal resolution is 170\,km. The horizontal
resolution perpendicular to the orbit track is 6\,km for all pressures.

\subsection*{Precision}

The estimated precision for SO\cs2 is \texttildelow3\,ppbv for all heights
between 215\,--\,10\,hPa. The precisions are set to negative values in
situations when the retrieved precision is larger than 50\% of the a
priori precision -- an indication that the data is biased toward the a
priori value.

\subsection*{Accuracy}

The values for accuracy are based on the approach used for the v2.2
systematic error analysis performed on the MLS measurement system
\citep{ReadEtal07}, updated for \thisv. The accuracy is estimated to be
\texttildelow5\,ppbv for pressures less than 147\,hPa increasing to
\texttildelow20\,ppbv at 215\,hPa.

\subsection*{Data screening}
\begin{description}
\item[Pressure range:] \screeningrule{215\,--\,10.0\,hPa} \standardrangestatement
\item[Estimated Precision:] \screeningrule{Values with negative
    precision can be used, though with caution.}

  Although it is generally recommended not to use values where precision
  is flagged negative, SO\cs2 is an exception and it is OK to use values
  with negatively flagged precision (provided that the entire profile is
  not so flagged). High retrieved values of SO\cs2 at the higher
  pressures (e.g. 215 and 147~hPa) will also have larger precision
  values which are sometimes large enough to trigger the ``too much \emph{a
    priori} influence'' flag. While a priori influence is present and
  the retrieved value is probably smaller than reality because the
  retrieval is being pulled towards the a priori value of zero, this
  does not detract from the fact that greatly enhanced SO\cs2 is being
  reported, reflecting the detection of a plume.
\evenstatusrule
\item[Quality:] \qualityrule{0.6}
\item[Convergence:] \convergencerule{1.8}
\end{description}

\subsection*{Artifacts}

The product is unvalidated. There is a tendency for the \thisv SO\cs2
product, as for CO, to have spikes in the presence of clouds.

\subsection*{Review of comparisons with other datasets}

MLS has successfully detected SO\cs2 from 23 eruptions since launch.
These include Manam (Papua, New Guinea -- 3 events), Anatahan (Mariana
Islands), Sierra Negra (Galapogos Island), Soufriere Hills (Montserrat,
West Indies), Tunguraua (Ecuador), Rabaul (Papua New Guinea),
Nyamuragira (Congo), Piton de la Fournaise (Reunion Island), Jebel
al-Tair (Yemen), Okmok (Alaska), Kasatochi (Alaska), Dalaffilla
(Ethiopia), Redoubt (Alaska), Sarychev (Kuril Islands, Russia), Pacaya
(Guatamala), Merapi (Indonesia), Grimsvoth (iceland), Puyehue-Cordon
Calle (Chile), Nabro (Eritrea) and Paluweh (Indonesia).

\begin{figure}
\begin{center}
\dontincludegraphics[width=0.7\linewidth]{figures/kosatochi_11082008}
\dontincludegraphics[width=0.7\linewidth]{figures/kosatochi_12082008}
\end{center}
\caption{An overlay of MLS measurement tracks on an OMI SO\cs2
  measurement on 11 and 12 August 2008 (separate maps) showing the
  dispersial of SO\cs2 from the Kasatochi eruption (8 Aug 2008, black
  triangle). The color scale indicates the SO\cs2 column measured by
  OMI\@. The daytime MLS tracks are small open circles and the nighttime
  tracks are filled black. When the calculated column from MLS exceeds 1
  DU, that measurement is indicated by a larger open circle filled with
  the color of the column measurement as indicated by the color scale
  below (same as for OMI)\@. The panels at right show all the measured
  profiles covering the area shown in the maps for SO\cs2 and
  HCl. Profiles where the MLS column calculation exceeds 1 DU are
  highlighted in red.}
\label{fig:so2overlay}
\end{figure}

Figure~\ref{fig:so2overlay} shows an overlay comparison of column SO\cs2
measured by OMI and the same calculated by MLS for two days following
the Kasatochi eruption. It is clear that MLS detects the main plume
dispersal features. It also appears that MLS columns are much smaller
than those from OMI. Interpreting the significance of this is not
straightforward given that OMI has to make assumptions regarding the
profile shape and can observe SO\cs2 down to the boundary layer. The MLS
column begins at 215~hPa and integrates upward neglecting the
tropospheric contributions. Another limitation is that OMI can only make
measurements during the day whereas MLS can make them day and night.
Since the plume is moving relatively quickly over the 12 hour
measurement separation time, MLS nighttime measurements often miss
and/or detect plume features differently than OMI\@. MLS vertically
resolved measurement shows that this plume has separated into distinct
layers at different altitudes.  The altitude of SO\cs2 measured by MLS
usually coincides with the altitude of aerosols measured by CALIPSO,
providing supporting evidence for the validity of the
vertically-resolved MLS observations.

\begin{table}
\caption{Summary of MLS \thisv SO\cs2 product.}
\label{tab:SO2Summary}
\vspace{\tablecaptionsep}
\begin{minipage}{\linewidth}
\centering\begin{tabular}{ccccl}
\toprule
\parbox[m]{6em}{\bfseries\centering Pressure / hPa} &
\parbox[m]{6em}{\bfseries\centering Resolution V~$\times$~H \ km} &
\parbox[m]{8em}{\bfseries\centering Single profile precision
\footnote{Absolute error in percent} / ppbv} &
\parbox[m]{6em}{\bfseries\centering Accuracy / ppbv} &
\bfseries Comments \\
\midrule
$<$ 10 & --- & --- & --- &  Unsuitable for scientific use \\
     10 & 3 $\times$ 180 &  3.5 &  6  & \\
     15 & 3 $\times$ 180 &  3.5 &  3  & \\
     22 & 3 $\times$ 180 &  3.2 &  4  & \\
     32 & 3 $\times$ 180 &  3.2 &  5  & \\
     46 & 3 $\times$ 180 &  3.0 &  5  & \\
     68 & 3 $\times$ 180 &  3.0 &  6  & \\
    100 & 3 $\times$ 180 &  3.0 &  6  & \\
    147 & 3 $\times$ 180 &  3.1 &  10 & \\
    215 & 3 $\times$ 180 &  3.8 &  20 & \\
$>$215 & --- & --- & --- & Unsuitable for scientific use \\
\bottomrule
\end{tabular}
\end{minipage}
\end{table}


%%% Local Variables:
%%% mode: latex
%%% TeX-master: "v4-x-report"
%%% End:
